% §5 The Practical Question: Which Cascade Architecture?
Practitioners currently choose between four concrete cascade configurations for a realtime,
tool-using, long-session assistant:
a serial graph orchestrator (no decoupling, no filler),
the same with manual filler,
the same with model-uttered (LLM) filler, and
a decoupled, event-driven runtime with tools on a background clock.
The two-floor model assigns each configuration its onset and cessation behaviour in one table
(Table~\ref{tab:architectures}).

\begin{table}[t]
\centering
\begin{tabular}{p{0.30\textwidth}p{0.32\textwidth}p{0.30\textwidth}}
\toprule
\textbf{Architecture} & \textbf{Onset} & \textbf{Cessation} \\
\midrule
Serial graph (all clocks coupled) & pays reducible latency serially &
can cancel generation; checkpoint preserves state, but the coupled loop drags tool-waits and
delivery into the stop path \\
\midrule
Serial graph + manual filler & onset masked at delivery layer &
manual filler shortens the interruption tail: the classic barge-in-friendly device \\
\midrule
Serial graph + LLM filler & onset masked but fragment may bleed or duplicate in the committed
text & no cessation gain; unretractable filler corrupts delivered content per
Corollary~\ref{cor:filler} \\
\midrule
Event-driven, decoupled clocks & hides reducible latency during user speech &
owns both floors: generation-interruptible, execution-persistent, atomically committed batches
(Lemma~\ref{lem:serializability}) \\
\bottomrule
\end{tabular}
\caption{The four cascade architectures mapped onto the two-floor model. The event-driven row
is a sufficiency statement (Definitions~\ref{def:interruptibility} and~\ref{def:clocks}), not
an in-principle uniqueness claim: a graph orchestrator that hosts async nodes, background
dispatch, and checkpoint reentry can emulate it at expressiveness cost.
\label{tab:architectures}}
\end{table}

Two consequences follow.
First, for single-hop turns there is no reducible latency to hide and no tail to manage: under
a fixed in-flight policy all four configurations meet the same floor, so orchestration choice
is noise.
Second, for tool and long-duration turns the configurations separate sharply: tool round-trips
land on the onset floor unless dispatched to a background clock
(Theorem~\ref{thm:tool_hiding}), and interruption behaviour is governed by the delivery tail
(Theorem~\ref{thm:cessation_floor}).
This separation is the pattern that tool-agent latency evaluations reveal: they measure
accuracy and liveness jointly but do not predict liveness from tool-DAG structure
\citep{taubench2024}, whereas our model predicts exactly the turns where structure dominates.
An assistant that mostly answers from memory can stay on a serial graph with manual filler;
an assistant whose turns habitually touch the bank, the card service, or multi-step skills
needs the decoupled runtime to avoid paying each tool round-trip inside the user's perceived
silence.
